% Perbandingan mode ECB dan CBC.
% Dirujuk dari section-02.tex dengan % Perbandingan mode ECB dan CBC.
% Dirujuk dari section-02.tex dengan % Perbandingan mode ECB dan CBC.
% Dirujuk dari section-02.tex dengan % Perbandingan mode ECB dan CBC.
% Dirujuk dari section-02.tex dengan \input{figures/mode-ecb-cbc}.
\begin{figure}[htbp]
    \centering
    \begin{tikzpicture}[
        kotak/.style={draw, rounded corners, align=center, font=\small, inner sep=3pt,
                      minimum width=0.95cm, minimum height=0.6cm},
        blok/.style={kotak, fill=blue!6},
        enc/.style={kotak, fill=orange!12, minimum width=1.1cm},
        cip/.style={kotak, fill=green!10},
        xor/.style={draw, circle, inner sep=1pt, font=\small, fill=yellow!18,
                    minimum size=0.5cm},
        panah/.style={-{Latex[length=2mm]}, thick},
    ]
        % ================= ECB =================
        \node[font=\small\bfseries] at (3,0.9) {Electronic Code Book (ECB)};
        \foreach \i/\x in {1/0, 2/3, 3/6} {
            \node[blok] (p\i) at (\x,0)   {$P_{\i}$};
            \node[enc]  (e\i) at (\x,-1.3) {$E_K$};
            \node[cip]  (c\i) at (\x,-2.6) {$C_{\i}$};
            \draw[panah] (p\i) -- (e\i);
            \draw[panah] (e\i) -- (c\i);
        }
        \node[font=\scriptsize, align=center] at (3,-3.5)
            {setiap blok dienkripsi sendiri-sendiri: blok plainteks yang sama\\ selalu menghasilkan cipherteks yang sama};

        % ================= CBC =================
        \node[font=\small\bfseries] at (3,-5.0) {Cipher Block Chaining (CBC)};
        \foreach \i/\x in {1/0, 2/3, 3/6} {
            \node[blok] (q\i) at (\x,-6.0) {$P_{\i}$};
            \node[xor]  (x\i) at (\x,-7.0) {$\oplus$};
            \node[enc]  (f\i) at (\x,-8.0) {$E_K$};
            \node[cip]  (d\i) at (\x,-9.0) {$C_{\i}$};
            \draw[panah] (q\i) -- (x\i);
            \draw[panah] (x\i) -- (f\i);
            \draw[panah] (f\i) -- (d\i);
        }
        \node[blok, fill=gray!12] (iv) at (-2.6,-7.0) {$IV$};
        \draw[panah] (iv) -- (x1);
        \draw[panah, dashed] (d1.east) -- node[sloped, above, font=\scriptsize] {$C_1$} (x2.west);
        \draw[panah, dashed] (d2.east) -- node[sloped, above, font=\scriptsize] {$C_2$} (x3.west);
        \node[font=\scriptsize, align=center] at (3,-9.9)
            {setiap blok di-XOR dengan cipherteks sebelumnya, sehingga blok yang sama\\ menghasilkan cipherteks berbeda sesuai posisinya};
    \end{tikzpicture}
    \caption{Perbandingan mode ECB yang independen dengan mode CBC yang merantai blok melalui operasi XOR}
    \label{fig:mode-ecb-cbc}
\end{figure}
.
\begin{figure}[htbp]
    \centering
    \begin{tikzpicture}[
        kotak/.style={draw, rounded corners, align=center, font=\small, inner sep=3pt,
                      minimum width=0.95cm, minimum height=0.6cm},
        blok/.style={kotak, fill=blue!6},
        enc/.style={kotak, fill=orange!12, minimum width=1.1cm},
        cip/.style={kotak, fill=green!10},
        xor/.style={draw, circle, inner sep=1pt, font=\small, fill=yellow!18,
                    minimum size=0.5cm},
        panah/.style={-{Latex[length=2mm]}, thick},
    ]
        % ================= ECB =================
        \node[font=\small\bfseries] at (3,0.9) {Electronic Code Book (ECB)};
        \foreach \i/\x in {1/0, 2/3, 3/6} {
            \node[blok] (p\i) at (\x,0)   {$P_{\i}$};
            \node[enc]  (e\i) at (\x,-1.3) {$E_K$};
            \node[cip]  (c\i) at (\x,-2.6) {$C_{\i}$};
            \draw[panah] (p\i) -- (e\i);
            \draw[panah] (e\i) -- (c\i);
        }
        \node[font=\scriptsize, align=center] at (3,-3.5)
            {setiap blok dienkripsi sendiri-sendiri: blok plainteks yang sama\\ selalu menghasilkan cipherteks yang sama};

        % ================= CBC =================
        \node[font=\small\bfseries] at (3,-5.0) {Cipher Block Chaining (CBC)};
        \foreach \i/\x in {1/0, 2/3, 3/6} {
            \node[blok] (q\i) at (\x,-6.0) {$P_{\i}$};
            \node[xor]  (x\i) at (\x,-7.0) {$\oplus$};
            \node[enc]  (f\i) at (\x,-8.0) {$E_K$};
            \node[cip]  (d\i) at (\x,-9.0) {$C_{\i}$};
            \draw[panah] (q\i) -- (x\i);
            \draw[panah] (x\i) -- (f\i);
            \draw[panah] (f\i) -- (d\i);
        }
        \node[blok, fill=gray!12] (iv) at (-2.6,-7.0) {$IV$};
        \draw[panah] (iv) -- (x1);
        \draw[panah, dashed] (d1.east) -- node[sloped, above, font=\scriptsize] {$C_1$} (x2.west);
        \draw[panah, dashed] (d2.east) -- node[sloped, above, font=\scriptsize] {$C_2$} (x3.west);
        \node[font=\scriptsize, align=center] at (3,-9.9)
            {setiap blok di-XOR dengan cipherteks sebelumnya, sehingga blok yang sama\\ menghasilkan cipherteks berbeda sesuai posisinya};
    \end{tikzpicture}
    \caption{Perbandingan mode ECB yang independen dengan mode CBC yang merantai blok melalui operasi XOR}
    \label{fig:mode-ecb-cbc}
\end{figure}
.
\begin{figure}[htbp]
    \centering
    \begin{tikzpicture}[
        kotak/.style={draw, rounded corners, align=center, font=\small, inner sep=3pt,
                      minimum width=0.95cm, minimum height=0.6cm},
        blok/.style={kotak, fill=blue!6},
        enc/.style={kotak, fill=orange!12, minimum width=1.1cm},
        cip/.style={kotak, fill=green!10},
        xor/.style={draw, circle, inner sep=1pt, font=\small, fill=yellow!18,
                    minimum size=0.5cm},
        panah/.style={-{Latex[length=2mm]}, thick},
    ]
        % ================= ECB =================
        \node[font=\small\bfseries] at (3,0.9) {Electronic Code Book (ECB)};
        \foreach \i/\x in {1/0, 2/3, 3/6} {
            \node[blok] (p\i) at (\x,0)   {$P_{\i}$};
            \node[enc]  (e\i) at (\x,-1.3) {$E_K$};
            \node[cip]  (c\i) at (\x,-2.6) {$C_{\i}$};
            \draw[panah] (p\i) -- (e\i);
            \draw[panah] (e\i) -- (c\i);
        }
        \node[font=\scriptsize, align=center] at (3,-3.5)
            {setiap blok dienkripsi sendiri-sendiri: blok plainteks yang sama\\ selalu menghasilkan cipherteks yang sama};

        % ================= CBC =================
        \node[font=\small\bfseries] at (3,-5.0) {Cipher Block Chaining (CBC)};
        \foreach \i/\x in {1/0, 2/3, 3/6} {
            \node[blok] (q\i) at (\x,-6.0) {$P_{\i}$};
            \node[xor]  (x\i) at (\x,-7.0) {$\oplus$};
            \node[enc]  (f\i) at (\x,-8.0) {$E_K$};
            \node[cip]  (d\i) at (\x,-9.0) {$C_{\i}$};
            \draw[panah] (q\i) -- (x\i);
            \draw[panah] (x\i) -- (f\i);
            \draw[panah] (f\i) -- (d\i);
        }
        \node[blok, fill=gray!12] (iv) at (-2.6,-7.0) {$IV$};
        \draw[panah] (iv) -- (x1);
        \draw[panah, dashed] (d1.east) -- node[sloped, above, font=\scriptsize] {$C_1$} (x2.west);
        \draw[panah, dashed] (d2.east) -- node[sloped, above, font=\scriptsize] {$C_2$} (x3.west);
        \node[font=\scriptsize, align=center] at (3,-9.9)
            {setiap blok di-XOR dengan cipherteks sebelumnya, sehingga blok yang sama\\ menghasilkan cipherteks berbeda sesuai posisinya};
    \end{tikzpicture}
    \caption{Perbandingan mode ECB yang independen dengan mode CBC yang merantai blok melalui operasi XOR}
    \label{fig:mode-ecb-cbc}
\end{figure}
.
\begin{figure}[htbp]
    \centering
    \begin{tikzpicture}[
        kotak/.style={draw, rounded corners, align=center, font=\small, inner sep=3pt,
                      minimum width=0.95cm, minimum height=0.6cm},
        blok/.style={kotak, fill=blue!6},
        enc/.style={kotak, fill=orange!12, minimum width=1.1cm},
        cip/.style={kotak, fill=green!10},
        xor/.style={draw, circle, inner sep=1pt, font=\small, fill=yellow!18,
                    minimum size=0.5cm},
        panah/.style={-{Latex[length=2mm]}, thick},
    ]
        % ================= ECB =================
        \node[font=\small\bfseries] at (3,0.9) {Electronic Code Book (ECB)};
        \foreach \i/\x in {1/0, 2/3, 3/6} {
            \node[blok] (p\i) at (\x,0)   {$P_{\i}$};
            \node[enc]  (e\i) at (\x,-1.3) {$E_K$};
            \node[cip]  (c\i) at (\x,-2.6) {$C_{\i}$};
            \draw[panah] (p\i) -- (e\i);
            \draw[panah] (e\i) -- (c\i);
        }
        \node[font=\scriptsize, align=center] at (3,-3.5)
            {setiap blok dienkripsi sendiri-sendiri: blok plainteks yang sama\\ selalu menghasilkan cipherteks yang sama};

        % ================= CBC =================
        \node[font=\small\bfseries] at (3,-5.0) {Cipher Block Chaining (CBC)};
        \foreach \i/\x in {1/0, 2/3, 3/6} {
            \node[blok] (q\i) at (\x,-6.0) {$P_{\i}$};
            \node[xor]  (x\i) at (\x,-7.0) {$\oplus$};
            \node[enc]  (f\i) at (\x,-8.0) {$E_K$};
            \node[cip]  (d\i) at (\x,-9.0) {$C_{\i}$};
            \draw[panah] (q\i) -- (x\i);
            \draw[panah] (x\i) -- (f\i);
            \draw[panah] (f\i) -- (d\i);
        }
        \node[blok, fill=gray!12] (iv) at (-2.6,-7.0) {$IV$};
        \draw[panah] (iv) -- (x1);
        \draw[panah, dashed] (d1.east) -- node[sloped, above, font=\scriptsize] {$C_1$} (x2.west);
        \draw[panah, dashed] (d2.east) -- node[sloped, above, font=\scriptsize] {$C_2$} (x3.west);
        \node[font=\scriptsize, align=center] at (3,-9.9)
            {setiap blok di-XOR dengan cipherteks sebelumnya, sehingga blok yang sama\\ menghasilkan cipherteks berbeda sesuai posisinya};
    \end{tikzpicture}
    \caption{Perbandingan mode ECB yang independen dengan mode CBC yang merantai blok melalui operasi XOR}
    \label{fig:mode-ecb-cbc}
\end{figure}
